\section{Anhang}


\subsection{Weboberfläche}
\label{app:Weboberfläche}

\subsubsection{UI-Entwurf}
\label{app:UI-Entwurf}
Zu sehen ist die Übersichtsseite der offenen Tasks. Neue Tasks werden fett-gedruckt dargestellt. Die Zeile auf der sich die Maus
befindet wird bereits hervorgehoben. Daher soll sich dies auch bei dem farbigen Balken, links jeden Eintrages, wiederspiegeln.
(Siehe sechster Eintrag von oben)
\begin{figure}[htb]
\centering
\includegraphicsKeepAspectRatio{Anhang/ui_entwurf.png}{1.0}
\caption{UI-Entwurf}
\end{figure}

\subsubsection{listOpenTasks.xhtml (Auszug)}
\label{app:listOpenTasks}
\lstinputlisting[language=HTML]{Listings/listOpenTasks.xhtml}

\clearpage
\subsubsection{main.less (Auszug)}
\label{app:main.less}
\lstinputlisting[language=CSS]{Listings/main.css}

\subsubsection{masterTemplate.xhtml (Auszug)}
\label{app:masterTemplate}
\lstinputlisting[language=HTML]{Listings/masterTemplate.xhtml}


\clearpage
\subsection{Tabellen}
\label{Tabellen}

\subsubsection{Laptop-Spezifikation}
\label{app:Laptop-Spezifikation}
\newTabelle{Laptop-Spezifikation}{tab:Laptop-Spezifikation}{Laptop.tex}

\subsubsection{Auszug der vorhandene Paketstruktur}
\label{app:Paketstruktur}
Ein vollständiges Beispielpacket wäre: \Code{com.<mandant>.contrax.contract.domain.model.transactions}
\newTabelle{Packetstruktur: Fachliche Saeulen}{tab:Fachliche Saeulen}{Saeulen.tex}
\newTabelle{Packetstruktur: Technische Layer}{tab:Technische Layer}{Layer.tex}

\clearpage
\subsection{Quellcode}
\label{app:Quellcode}
Bei Java-Dateien wurden zur Übersichtlichkeit sämtliche Importierungen mit "`\Code{import *;}"' abgekürzt. Ebent so sämtliche
Kommentare, einschließlich JavaDoc, sofern diese nicht maßgeblich zum Verständnis beitragen. Es wird alphabetisch nach Dateinamen
aufgelistet.

\subsubsection{AbstractNotificationService.java}
\label{app:AbstractNotificationService}
\lstinputlisting[style=Java]{Listings/AbstractNotificationService.java}

\subsubsection{AtmosphereConfig.java}
\label{app:AtmosphereConfig}
\lstinputlisting[style=Java]{Listings/AtmosphereConfig.java}

\subsubsection{main.js (Auszug)}
\label{app:main.js}
Die Methode, welche den Hover-Effekt in der Task-Übersicht ermöglicht, und die Methode \\
\Methode{configureListeners(request,socket)} wurden aufgrund der Anhangsbegrenzung ausgelassen. In Letzterer werden \sog Listener,
wie \Methode{onOpen} und \Methode{onTransportFailure} definiert. Der \Methode{onMessage}-Listener wurde wegen seiner Wichtigkeit
gesondert implementiert.
\lstinputlisting[language=JavaScript]{Listings/main.js}

\subsubsection{NotificationController.java}
\label{app:NotificationController}
\lstinputlisting[style=Java]{Listings/NotificationController.java}

\subsubsection{NotificationServiceHandler.java}
\label{app:NotificationServiceHandler}
\lstinputlisting[style=Java]{Listings/NotificationServiceHandler.java}

\subsubsection{OpenTaskListController.java}
\label{app:OpenTaskListController}
\lstinputlisting[style=Java]{Listings/OpenTaskListController.java}

\subsubsection{pom.xml (Auszug)}
\label{app:pom.xml}
\lstinputlisting[style=XML]{Listings/pompom.xml}

\clearpage
\subsubsection{PostRenderView-Event}
\label{app:PostRenderView-Event}
Zum leztlichen Eintragen des Events in den JSF-Lebenszyklus, muss der \sog "`Phaselistener"' als dieser in die
\Code{faces-config.xml} eingetragen werden.
\lstinputlisting[style=Java]{Listings/PostRenderViewEvent.java}
\lstinputlisting[style=Java]{Listings/PostRenderViewListener.java}

\subsubsection{TaskNotificationEvent.java}
\label{app:TaskNotificationEvent}
\lstinputlisting[style=Java]{Listings/TaskNotificationEvent.java}

\subsubsection{TaskNotificationService.java}
\label{app:TaskNotificationService}
\lstinputlisting[style=Java]{Listings/TaskNotificationService.java}

\subsubsection{TaskQueryServiceImpl.java (Auszug)}
\label{app:TaskQueryServiceImpl}
Interface:
\lstinputlisting[style=Java]{Listings/TaskQueryService.java}
Implementierung:
\lstinputlisting[style=Java]{Listings/TaskQueryServiceImpl.java}

\subsubsection{UserStatus(Repository).java}
\label{app:UserStatusRepository}
\lstinputlisting[style=Java]{Listings/UserStatus.java}
\lstinputlisting[style=Java]{Listings/UserStatusRepository.java}

\clearpage
\subsubsection{VmsWebAppInitializer.java (Auszug)}
\label{app:VmsWebAppInitializer}
\lstinputlisting[style=Java]{Listings/VmsWebAppInitializer.java}


\clearpage
\subsection{Weitere Grafiken}
\label{app:Grafiken}
Alle Grafiken sind in ihrer Größe reduziert. In den Klassendiagrammen werden ausschließlich Klassen betrachtet die
hinzugefügt oder inhaltlich geändert wurden. Vererbungen und Implementierungen sind, sofern von Relevanz, angegeben.

\subsubsection{UML-Klassendiagramm (Übersicht)}
\label{app:UML-Klassendiagramm (Übersicht)}
\begin{figure}[htb]
\centering
\includegraphicsKeepAspectRatio{Anhang/uml-diagram.png}{1}
\caption{UML-Klassendiagramm (Übersicht)}
\end{figure}

\clearpage
\subsubsection{UML-Klassendiagramm (Komplex)}
\label{app:UML-Klassendiagramm (Komplex)}
Eine Blau-geschriebene Methode weist darauf hin, dass diese verändert wurde. Eine in Grün geschrieben, dass diese neu hinzugefügt
wurde. Beinhaltet eine dargestellte Klasse nur Letzteres ist davon auszugehen, dass diese gänzlich Neu ist.
\begin{figure}[htb]
\centering
\includegraphicsKeepAspectRatio{Anhang/uml-diagram-full3.png}{1}
\caption{UML-Klassendiagramm (Komplex)}
\end{figure}





